% Generated from locale/tr/content/incompleteness/representability-in-q/prim-rec.tex; do not hand-edit.


\olfileid[tr]{inc}{req}{pri}
\olsection{İlkel Özyinelemenin Benzetilmesi}

Artık beta fonksiyonunu kullanarak ilkel özyinelemeyle tanımın düzenli
minimizasyonla ``benzetilebildiğini'' gösterebiliriz. $f(\vec x)$ ve
$g(\vec x,y,z)$ fonksiyonlarımız olduğunu varsayın. O zaman $f$ ile $g$'den
ilkel özyinelemeyle tanımlanan $h(x,\vec z)$ fonksiyonu şöyledir:
\begin{align*}
h(\vec x, 0) & =  f(\vec x) \\
h(\vec x, y+1) & =  g(\vec x, y, h(\vec x, y)).
\end{align*}
$h$'nin yalnızca bileşke ve düzenli minimizasyon kullanılarak $f$ ile
$g$'den, temel fonksiyonlardan ve bunlardan bileşke ile düzenli minimizasyon
kullanılarak tanımlanan fonksiyonlardan ($\beta$ gibi) tanımlanabildiğini
göstermeliyiz.

\begin{lem}
\ollabel{lem:prim-rec}
$h$, $f$ ile $g$'den ilkel özyineleme kullanılarak tanımlanabiliyorsa;
$f$, $g$, $\Zero$, $\Succ$, $\Proj{n}{i}$, $\Add$, $\Mult$ ve $\Char{=}$
fonksiyonlarından bileşke ve düzenli minimizasyon kullanılarak da
tanımlanabilir.
\end{lem}

\begin{proof}
Önce en küçük $d$ sayısını döndüren yardımcı bir $\hat h(\vec x,y)$
fonksiyonu tanımlayın; bu $d$, aşağıdakileri sağlayan bir diziyi kodlasın:
\begin{enumerate}
\item $(d)_0 = f(\vec x)$ ve
\item her $i<y$ için $(d)_{i+1}=g(\vec x,i,(d)_i)$.
\end{enumerate}
Burada artık $(d)_i$, $\beta(d,i)$'nin kısaltmasıdır. Başka bir deyişle
$\hat h$, $\tuple{h(\vec x,0),h(\vec x,1),\dots,h(\vec x,y)}$ dizisini
döndürür. $\hat h$'yi şöyle yazabiliriz:
\[
\hat h(\vec x, y) = \umin{d}{(\beta(d,0) = f(\vec x) \land \bforall{i <
  y}{\beta(d,i+1) = g(\vec x, i,\beta(d,i)})}.
\]
Dikkat: burada ilkel özyineleme gerekmez; yalnızca minimizasyon kullanılır.
Minimize ettiğimiz fonksiyon, \olref[bet]{lem:beta} beta fonksiyonu lemması
uyarınca düzenlidir.

Fakat artık
\[
h(\vec x, y) = \beta(\hat h(\vec x, y), y),
\]
olur; dolayısıyla $h$ yalnızca bileşke ve düzenli minimizasyon kullanılarak
temel fonksiyonlardan tanımlanabilir.
\end{proof}

